% =============================================================================
%  lec13.tex —— 第 13 讲 可表现 ∞-范畴（底本 PDF 第 104–109 页）
% =============================================================================
\lecture{13}{可表现 ∞-范畴}
\begin{center}\small\color{black!55}（底本 PDF 第 104–109 页）\end{center}

本讲引入可表现 $\infty$-范畴。

\section{Ind-完备化}

\nitem\label{L13-1.1} 设 $\cC$ 是本质小的 $\infty$-范畴，
\[
\iota\colon\cC\to\mathrm{PShv}(\cC)\defeq\Fun(\cC^{\op},\Grpd_\infty)
\]
是 Yoneda 嵌入。回忆：可以把 $\mathrm{PShv}(\cC)$ 视为在 $\cC$ 上自由添加小余极限所得的 $\infty$-范畴。本节考虑在 $\cC$ 上自由添加小滤过余极限所得的 $\infty$-范畴。

\begin{definition}\label{L13-1.2}
设 $\cC$ 是小 $\infty$-范畴。定义 $\Ind(\cC)\subset\mathrm{PShv}(\cC)$ 为 $\mathrm{PShv}(\cC)$ 中由 $\iota(\cC)\subset\mathrm{PShv}(\cC)$ 在小滤过余极限下生成的全子 $\infty$-范畴。称其为 $\cC$ 的\textbf{Ind-完备化}。$\Ind(\cC)$ 中的对象称为 $\cC$ 中的\textbf{ind-对象}。

我们滥用记号，把所得的满忠实函子
\[
\iota\colon\cC\to\Ind(\cC)
\]
仍称为 Yoneda 嵌入。
\end{definition}

\yizhu{“Ind-完备化”与“ind-对象”保留了拉丁词干 ind（inductive 的缩写），因为该词已与“归纳（极限）”脱钩：Ind-完备化是自由添加小滤过余极限，而非添加归纳极限。相应地，对偶构造记作 Pro-完备化。\textbf{[来源：B]}}

\nitem\label{L13-1.3} 按定义，$\Ind(\cC)$ 中的任意对象都可以非唯一地写成
\[
\colim_{i\in I}\iota(x_i),
\]
其中 $I\to\cC$，$i\mapsto x_i$ 是 $\cC$ 中的小滤过图。我们常用记号
\[
\text{“$\colim_{i\in I}x_i$”}\defeq\colim_{i\in I}\iota(x_i)
\]
来表示这样的对象。

\begin{exercise}\label{L13-1.4}
证明
\[
\mathrm{Maps}_{\Ind(\cC)}\Big(\text{“$\colim_{i\in I}x_i$”},\text{“$\colim_{j\in J}y_j$”}\Big)\simeq\lim_{i\in I}\colim_{j\in J}\mathrm{Maps}_{\cC}(x_i,y_j).
\]
\end{exercise}

\nitem\label{L13-1.5} 回忆：$\mathrm{PShv}(\cC)$ 中的任意对象 $\mathcal{F}$ 都可以写成 $\iota(\cC)$ 中一个显式图的余极限：
\[
\mathcal{F}\simeq\colim_{x\in\cC_{/\mathcal{F}}}\iota(x),
\]
其中
\[
\cC_{/\mathcal{F}}\defeq\cC\times_{\mathrm{PShv}(\cC)}\mathrm{PShv}(\cC)_{/\mathcal{F}}
\]
是分类“$x\in\cC$ 连同态射 $\iota(x)\to\mathcal{F}$”的 $\infty$-范畴。

\begin{proposition}[HTT.5.3.5.4]\label{L13-1.6}
设 $\cC$ 是本质小的 $\infty$-范畴，$\mathcal{F}\in\mathrm{PShv}(\cC)$。则下述条件等价：
\begin{enumerate}[label=(\roman*),leftmargin=2.4em,itemsep=0.1em]
\item 对象 $\mathcal{F}$ 含于 $\Ind(\cC)$；
\item $\infty$-范畴 $\cC_{/\mathcal{F}}$ 是滤过的。
\end{enumerate}
此外，若 $\cC^{\op}$ 允许有限极限，则上述条件还等价于
\begin{enumerate}[label=(\roman*),leftmargin=2.4em,itemsep=0.1em,start=3]
\item 函子 $\mathcal{F}\colon\cC^{\op}\to\Grpd_\infty$ 保持有限极限。
\end{enumerate}
\end{proposition}

\begin{proposition}\label{L13-1.7}
设 $\cC$ 是本质小的 $\infty$-范畴。则 Yoneda 嵌入 $\iota\colon\cC\to\Ind(\cC)$ 保持并检测有限余极限。
\end{proposition}

\begin{proof}[证明梗概]
设 $I\to\cC$，$i\mapsto x_i$ 是允许余极限的有限图。则对任意对象“$\colim_{j\in J}y_j$”$\in\Ind(\cC)$，有
\[
\begin{aligned}
&\mathrm{Maps}_{\Ind(\cC)}\Big(\iota\big(\colim_{i\in I}x_i\big),\text{“$\colim_{j\in J}y_j$”}\Big)\\
&\qquad\simeq\colim_{j\in J}\mathrm{Maps}_{\cC}\Big(\colim_{i\in I}x_i,y_j\Big)\\
&\qquad\simeq\colim_{j\in J}\lim_{i\in I}\mathrm{Maps}_{\cC}(x_i,y_j)\\
&\qquad\simeq\lim_{i\in I}\colim_{j\in J}\mathrm{Maps}_{\cC}(x_i,y_j)\\
&\qquad\simeq\lim_{i\in I}\mathrm{Maps}_{\Ind(\cC)}\Big(\iota(x_i),\text{“$\colim_{j\in J}y_j$”}\Big)\\
&\qquad\simeq\mathrm{Maps}_{\Ind(\cC)}\Big(\colim_{i\in I}\iota(x_i),\text{“$\colim_{j\in J}y_j$”}\Big).
\end{aligned}
\]
其中第一个与第四个等价来自练习~\ref{L13-1.4}。第二个与第五个等价来自可表函子保持极限这一事实。第三个等价来自滤过余极限与有限极限交换这一事实。
\end{proof}

\begin{exercise}\label{L13-1.8}
证明 $\iota\colon\cC\to\Ind(\cC)$ 保持并检测所有极限。
\end{exercise}

\nitem\label{L13-1.9} 回忆：$\infty$-范畴 $\mathrm{PShv}(\cC)$ 是 $\cC$ 的余完备化，意即它由下述泛性质刻画：
\[
\mathrm{LFun}(\mathrm{PShv}(\cC),\cD)\xrightarrow{\ \simeq\ }\Fun(\cC,\cD),
\]
其中左端是保持小余极限的函子 $\mathrm{PShv}(\cC)\to\cD$ 构成的 $\infty$-范畴。$\infty$-范畴 $\Ind(\cC)$ 也可以由类似的泛性质刻画，只需把小余极限换成小滤过余极限。

\begin{definition}\label{L13-1.10}
设 $\cD$ 是允许小滤过余极限的 $\infty$-范畴。若函子 $F\colon\cD\to\cD'$ 保持小滤过余极限，则称 $F$ 是\textbf{连续的}。令
\[
\Fun(\cD,\cD')_{\mathrm{cont}}\subset\Fun(\cD,\cD')
\]
为连续函子构成的全子 $\infty$-范畴。
\end{definition}

\yizhu{此处的“连续”是范畴论意义下的连续，即保持（小滤过）余极限，与拓扑学中的连续性无关；中文文献亦作“余极限连续”。\textbf{[来源：B]}}

\begin{proposition}[HTT.5.3.5.10]\label{L13-1.11}
设 $\cC$ 是小 $\infty$-范畴，$\cD$ 是允许小滤过余极限的 $\infty$-范畴。则 Yoneda 嵌入 $\iota\colon\cC\to\Ind(\cC)$ 诱导等价
\[
\Fun(\Ind(\cC),\cD)_{\mathrm{cont}}\xrightarrow{\ -\circ\iota\ }\Fun(\cC,\cD),
\]
其逆由
\[
\Fun(\cC,\cD)\to\Fun(\Ind(\cC),\cD)_{\mathrm{cont}},\qquad F\mapsto\mathrm{LKE}_\iota F
\]
给出。
\end{proposition}

\section{紧对象}

\begin{definition}\label{L13-2.1}
设 $\cD$ 是允许小滤过余极限的 $\infty$-范畴。若函子\footnote{由于我们并不假设 $\cD$ 局部小，函子 $\mathrm{Maps}_{\cD}(x,-)$ 的靶应当是 $\widehat{\Grpd}_\infty$，它包含不必小的 $\infty$-群胚。} $\mathrm{Maps}_{\cD}(x,-)$ 连续，即保持小滤过余极限，则称对象 $x\in\cD$ 是\textbf{紧的}。令 $\cD^{\mathrm{cpt}}\subset\cD$ 为 $\cD$ 中紧对象构成的全子 $\infty$-范畴。
\end{definition}

\begin{proposition}[HTT.5.4.2.4]\label{L13-2.2}
设 $\cC$ 是小 $\infty$-范畴。则 $\Ind(\cC)^{\mathrm{cpt}}$ 是 $\Ind(\cC)$ 的满足下述两条的最小全子 $\infty$-范畴：
\begin{enumerate}[label=(\roman*),leftmargin=2.4em,itemsep=0.1em]
\item 它包含 $\iota(\cC)$；
\item 它是幂等完备的\footnote{这是指 $\Ind(\cC)^{\mathrm{cpt}}$ 允许 $\mathrm{Idem}$-指标的余极限。}。
\end{enumerate}
\end{proposition}

\begin{exercise}\label{L13-2.3}
请自行证明 $\Ind(\cC)^{\mathrm{cpt}}$ 满足这两条性质。
\end{exercise}

\begin{exercise}\label{L13-2.4}
证明 $\Ind(\cC)$ 是幂等完备的。由此推出：(ii) 可以换成
\begin{enumerate}[label=(\roman*'),leftmargin=2.8em,itemsep=0.1em,start=2]
\item 嵌入 $\cC'\subset\Ind(\cC)$ 在收缩下封闭\footnote{这是指：若 $z\in\cC'$ 且 $y\in\Ind(\cC)$ 使得 $\id_y$ 分解为 $y\to z\to y$，则 $y\in\cC'$。}。
\end{enumerate}
\end{exercise}

\begin{definition}\label{L13-2.5}
设 $\cD$ 是 $\infty$-范畴。称 $\cD$ 是\textbf{紧生成的}，若
\begin{enumerate}[label=(\alph*),leftmargin=2.4em,itemsep=0.1em]
\item 存在本质小的 $\infty$-范畴 $\cC$，使得 $\cD\simeq\Ind(\cC)$；
\item $\infty$-范畴 $\cD$ 允许所有小余极限。
\end{enumerate}
\end{definition}

\begin{proposition}[HTT.5.5.1.1]\label{L13-2.6}
设 $\cD$ 是 $\infty$-范畴。则下述条件等价：
\begin{itemize}[leftmargin=2.2em,itemsep=0.1em]
\item $\infty$-范畴 $\cD$ 是紧生成的；
\item $\infty$-范畴 $\cD^{\mathrm{cpt}}$ 允许有限余极限，并在小滤过余极限下生成 $\cD$；
\item 存在允许有限余极限的小 $\infty$-范畴 $\cC$，使得 $\cD\simeq\Ind(\cC)$。
\end{itemize}
\end{proposition}

\begin{exercise}\label{L13-2.7}
证明：若 $\cD$ 是紧生成的，则 $\cD$ 中的滤过余极限与有限极限交换。
\end{exercise}

\begin{exercise}\label{L13-2.8}
设 $F\colon\cC\rightleftarrows\cD\colon G$ 是伴随，且 $\cC$ 紧生成。证明：$G$ 连续当且仅当 $F$ 保持紧对象。
\end{exercise}

\section{可表现 $\infty$-范畴}

\nitem\label{L13-3.1} 回忆：$\infty$-范畴 $J$ 是滤过的，是指任意有限图 $K\to J$ 都能延拓为 $K^{\triangleright}\to J$。若把上述“任意有限图”换成“任意 $\kappa$-小图”，其中 $\kappa$ 是正则基数\footnote{直观地说，基数用来度量集合的大小。给定基数 $\kappa$，若集合 $S$ 的大小严格小于 $\kappa$，则称 $S$ 是 $\kappa$-小的。基数 $\kappa$ 是正则的，是指 $\kappa$-小集合的 $\kappa$-小并仍是 $\kappa$-小的。}，则前面关于滤过余极限的全部讨论仍然成立。例如，可以定义 $\kappa$-滤过图、$\kappa$-Ind-完备化 $\Ind_\kappa(-)$、$\kappa$-紧对象、$\kappa$-紧生成……相应结果的证明与 $\kappa$ 是可数基数 $\omega$ 的情形基本相同。

\begin{exercise}\label{L13-3.2}
设 $\kappa<\kappa'$ 是正则基数。证明：
\begin{center}
\begin{tabular}{c}
$\kappa$-小 $\Rightarrow$ $\kappa'$-小\\
$\kappa$-滤过 $\Leftarrow$ $\kappa'$-滤过\\
$\kappa$-连续 $\Leftarrow$ $\kappa'$-连续\\
$\kappa$-紧 $\Rightarrow$ $\kappa'$-紧
\end{tabular}
\end{center}
\end{exercise}

\begin{definition}\label{L13-3.3}
设 $\kappa$ 是正则基数。若存在本质小的 $\infty$-范畴 $\cC$，使得 $\cD\simeq\Ind_\kappa(\cC)$，则称 $\infty$-范畴 $\cD$ 是\textbf{$\kappa$-可及的}。若 $\cD$ 对某个 $\kappa$ 是 $\kappa$-可及的，则称 $\cD$ 是\textbf{可及的}。
\end{definition}

\begin{definition}\label{L13-3.4}
设 $\cD$ 可及。若函子 $\cD\to\cD'$ 对某个正则基数 $\kappa$ 是 $\kappa$-连续的，则称它是\textbf{可及的}。
\end{definition}

\begin{lemma}[HTT.5.4.2.8]\label{L13-3.5}
设 $\cD$ 是 $\kappa$-可及的，则对 $\kappa'\gg\kappa$\footnote{关于 $\gg$ 的含义，见 HTT.A.2.6.3。}，$\cD$ 也是 $\kappa'$-可及的。
\end{lemma}

\begin{definition}\label{L13-3.6}
若 $\infty$-范畴 $\cD$ 可及且允许小余极限，则称 $\cD$ 是\textbf{可表现的}。
\end{definition}

\nitem\label{L13-3.7} 按定义，$\cD$ 可表现，当且仅当它对某个正则基数 $\kappa$ 是 $\kappa$-紧生成的。

\begin{theorem}[HTT.A.3.7.6]\label{L13-3.8}
$\infty$-范畴 $\cD$ 可表现，当且仅当存在组合单纯模型范畴 $\mathfrak{A}$，使得 $\cD\simeq\mathfrak{N}_{\bullet}(\mathfrak{A}^{\circ})$。
\end{theorem}

\yizhu{原书此处 $\mathfrak{N}_{\bullet}$ 指单纯脉（simplicial nerve），$\mathfrak{A}^{\circ}$ 指模型范畴 $\mathfrak{A}$ 的以全体上纤维对象为对象、以全体态射为态射的单纯子范畴。\textbf{[来源：B]}}

\begin{corollary}\label{L13-3.9}
$\infty$-范畴 $\Grpd_\infty$ 与 $\Cat_\infty$ 都是可表现的。
\end{corollary}

\begin{theorem}[HTT.5.5.2.4]\label{L13-3.10}
可表现的 $\infty$-范畴也允许小极限。
\end{theorem}

\begin{warning}\label{L13-3.11}
若 $\cD$ 可表现，则 $\cD^{\op}$ 几乎绝不是可及的。挑战：证明 $\Grpd^{\op}_\infty$ 不是可及的。
\end{warning}

\begin{theorem}[HTT.5.5.1.1]\label{L13-3.12}
$\infty$-范畴 $\cD$ 可表现，当且仅当存在本质小的 $\infty$-范畴 $\cC$ 与伴随
\[
F\colon\mathrm{PShv}(\cC)\rightleftarrows\cD\colon G,
\]
使得 $G$ 满忠实且 $G\circ F$ 可及。
\end{theorem}

\begin{remark}\label{L13-3.13}
上述定理是证明关于可表现 $\infty$-范畴的结论的有力工具。即：可以先在 $\cD=\mathrm{PShv}(\cC)$ 的情形证明，再形式地推出一般情形。
\end{remark}

\begin{proposition}[HTT.5.5.3.6]\label{L13-3.14}
设 $\cC$ 可表现，$K$ 是小单纯集。则 $\Fun(K,\cC)$ 可表现。
\end{proposition}

\begin{proposition}[HTT.5.5.3.10, 5.5.3.11]\label{L13-3.15}
设 $\cC$ 可表现，$u\colon K\to\cC$ 是小图。则 $\cC_{/u}$ 与 $\cC_{u/}$ 都可表现。
\end{proposition}

\section{可表性定理与伴随函子定理}

\begin{exercise}\label{L13-4.1}
可表现 $\infty$-范畴是局部小的，即 $\mathrm{Maps}(-,-)$ 本质小。
\end{exercise}

\begin{theorem}[HTT.5.5.2.2, 5.5.2.7]\label{L13-4.2}
设 $\cC$ 是可表现 $\infty$-范畴。
\begin{enumerate}[label=(\arabic*),leftmargin=2.4em,itemsep=0.1em]
\item 函子 $\cC^{\op}\to\Grpd_\infty$ 可表示，当且仅当它保持小极限。
\item 函子 $\cC\to\Grpd_\infty$ 余可表示，当且仅当它可及且保持小极限。
\end{enumerate}
\end{theorem}

\begin{remark}\label{L13-4.3}
由警告~\ref{L13-3.11}，要求 $\cC^{\op}\to\Grpd_\infty$ 可及是没有意义的。
\end{remark}

\begin{corollary}\label{L13-4.4}
设 $F\colon\cC\to\cD$ 是可表现 $\infty$-范畴之间的函子。
\begin{enumerate}[label=(\arabic*),leftmargin=2.4em,itemsep=0.1em]
\item 函子 $F$ 有右伴随，当且仅当它保持小余极限。
\item 函子 $F$ 有左伴随，当且仅当它可及且保持小极限。
\end{enumerate}
\end{corollary}

\yizhu{原文用 right adjointable、left adjointable 两个形容词，字面为“右/左可伴随的”，意即该函子有右/左伴随。译文直接写“有右伴随/有左伴随”，以免与“右伴随函子”混淆。\textbf{[来源：I]}}

\begin{proof}[证明梗概]
我们证明 (1)，把 (2) 留作练习。$F$ 有右伴随，当且仅当对任意 $y$，$\mathrm{Maps}(F(-),y)$ 可表示。由定理~\ref{L13-4.2}，这等价于：对任意 $y$，$\mathrm{Maps}(F(-),y)$ 保持小极限。由于可表函子保持并检测所有极限，这又等价于 $F(-)$ 保持小余极限。
\end{proof}

\section{可表现 $\infty$-范畴之间的左、右函子}

\begin{definition}\label{L13-5.1}
设
\[
\mathrm{Pr}^{L},\ \mathrm{Pr}^{R}
\]
是可表现 $\infty$-范畴构成的 $\infty$-范畴，其态射由左伴随（相应地，右伴随）函子给出。换言之，
\[
\mathrm{Maps}_{\mathrm{Pr}^{L}}(-,-)\defeq\mathrm{LFun}(-,-)^{\simeq},\qquad
\mathrm{Maps}_{\mathrm{Pr}^{R}}(-,-)\defeq\mathrm{RFun}(-,-)^{\simeq}.
\]
\end{definition}

\begin{proposition}[HTT.5.5.3.8]\label{L13-5.2}
设 $\cC$ 与 $\cD$ 是可表现 $\infty$-范畴。则 $\mathrm{LFun}(\cC,\cD)$ 可表现。
\end{proposition}

\begin{warning}\label{L13-5.3}
$\infty$-范畴 $\mathrm{RFun}(\cC,\cD)$ 几乎绝不是可表现的，因为
\[
\mathrm{RFun}(\cC,\cD)\simeq\mathrm{LFun}(\cD,\cC)^{\op}.
\]
\end{warning}

\begin{exercise}\label{L13-5.4}
$\mathrm{LFun}(\Grpd_\infty,\cD)$ 与 $\mathrm{RFun}(\Grpd^{\op}_\infty,\cD)$ 各是什么？$\mathrm{RFun}(\cD^{\op},\Grpd_\infty)$ 又如何？
\end{exercise}

\begin{exercise}\label{L13-5.5}
挑战：对任意可表现 $\infty$-范畴 $\cC$ 与 $\cD$，构造典范等价
\[
\mathrm{RFun}(\cC^{\op},\cD)\simeq\mathrm{RFun}(\cD^{\op},\cC).
\]
\end{exercise}

\begin{theorem}[HTT.5.5.3.4]\label{L13-5.6}
存在典范等价
\[
\mathrm{Pr}^{L}\simeq(\mathrm{Pr}^{R})^{\op},
\]
使得
\begin{itemize}[leftmargin=2.2em,itemsep=0.1em]
\item 它把可表现 $\infty$-范畴 $\cD$（视为 $\mathrm{Pr}^{L}$ 中的对象）送到 $\cD$（视为 $\mathrm{Pr}^{R}$ 中的对象）的一个等价对象；
\item 它把函子 $F\in\mathrm{LFun}(\cC,\cD)$ 送到它的一个右伴随 $F^{R}$。
\end{itemize}
\end{theorem}

\begin{remark}\label{L13-5.7}
特别地，给定图 $u\colon K\to\mathrm{Pr}^{L}$，我们可以取右伴随函子，得到图 $K^{\op}\to\mathrm{Pr}^{R}$。
\end{remark}

\begin{theorem}[HTT.5.5.3.13, 5.5.3.18]\label{L13-5.8}
$\infty$-范畴 $\mathrm{Pr}^{L}$ 与 $\mathrm{Pr}^{R}$ 都允许小余极限与小极限。而且遗忘函子\footnote{注意：可表现 $\infty$-范畴一般不是本质小的。因此我们需要用 $\widehat{\Cat}_\infty$ 而不是 $\Cat_\infty$。}
\[
\mathrm{Pr}^{L}\to\widehat{\Cat}_\infty,\qquad \mathrm{Pr}^{R}\to\widehat{\Cat}_\infty
\]
保持小极限。
\end{theorem}

\begin{remark}\label{L13-5.9}
回忆：$\widehat{\Cat}_\infty$ 中的极限是可计算的（第 9 讲注 2.9）。由此推出 $\mathrm{Pr}^{L}$ 与 $\mathrm{Pr}^{R}$ 中的小极限可计算。利用定理~\ref{L13-5.6}，$\mathrm{Pr}^{L}$ 与 $\mathrm{Pr}^{R}$ 中的小余极限可通过取伴随函子来计算。
\end{remark}

\appendixsection{紧空间}

\begin{exercise}\label{L13-A.1}
证明 $\Grpd_\infty$ 是紧生成的，且任意本质有限的 $\infty$-群胚都是紧对象。
\end{exercise}

\begin{exercise}\label{L13-A.2}
证明 $\Grpd_\infty$ 中的紧对象未必本质有限。
\end{exercise}

\nitem\label{L13-A.3} \textbf{建议阅读}\quad [Wal65] 与 [Wal66]。

\appendixsection{自由添加余极限}

\nitem\label{L13-B.1} 更一般地，对任意 $\infty$-范畴 $\cC$，都可以自由添加某几类余极限，同时保留 $\cC$ 中已存在的某一族余极限。

\begin{exercise}\label{L13-B.2}
给出刻画上述构造的泛性质。
\end{exercise}

\begin{example}\label{L13-B.3}
对本质小的 $\infty$-范畴 $\cC$，可以自由添加小筛余极限，得到所谓 $\cC$ 的\textbf{非交换导出 $\infty$-范畴}。
\end{example}

\nitem\label{L13-B.4} \textbf{建议阅读}\quad HTT.5.3.6、HTT.5.5.8 与 HTT.5.5.9。

\subsection*{参考文献}
\par\noindent\hangindent=5.2em\hangafter=1 [Wal65]\hspace{0.6em}Charles Terence Clegg Wall. Finiteness conditions for cw-complexes. \emph{Annals of Mathematics}, 81(1):56--69, 1965.
\par\noindent\hangindent=5.2em\hangafter=1 [Wal66]\hspace{0.6em}C Terence C Wall. Finiteness conditions for cw complexes. ii. \emph{Proceedings of the Royal Society of London. Series A, Mathematical and Physical Sciences}, pages 129--139, 1966.
